\documentclass{article}
\usepackage{amsmath,amssymb}
\usepackage[ruled,vlined,linesnumbered]{algorithm2e}

\DeclareMathOperator{\unit}{unit}
\SetKw{KwInParallel}{in parallel}

\begin{document}

\begin{algorithm}[t]
\caption{Communication layer: a hop-count gradient toward the goal sets each particle's polarity and curvity}
\label{alg:comm}
\KwIn{particle positions $\mathbf{x}_i$ and headings $\hat{\mathbf{e}}_i$; previous scores $s_i^{(t-1)}$;
      goal $\mathbf{g}$; view range $r$}
\KwOut{polarity $\mathbf{p}_i$, score $s_i^{(t)}$ and curvity $\kappa_i$ for every particle $i$}
\BlankLine
\For{each particle $i$}{
  $\mathcal{V}_i \leftarrow$ particles within range $r$ of $i$ with no wall in between\;
  \BlankLine
  \uIf{$i$ sees the goal}{
    $\mathbf{p}_i \leftarrow \unit(\mathbf{g} - \mathbf{x}_i)$\;
    $s_i^{(t)} \leftarrow 0$\;
  }
  \uElseIf{$\mathcal{V}_i = \emptyset$}{
    $\mathbf{p}_i \leftarrow \mathbf{0}$\;
    $s_i^{(t)} \leftarrow \infty$\;
  }
  \Else{
    $s^{*} \leftarrow \min_{j \in \mathcal{V}_i} s_j^{(t-1)}$\;
    $\mathcal{B}_i \leftarrow \{\, j \in \mathcal{V}_i : s_j^{(t-1)} = s^{*} \,\}$\;
    $\mathbf{p}_i \leftarrow \unit\bigl(\operatorname{mean}_{j \in \mathcal{B}_i} \mathbf{x}_j - \mathbf{x}_i\bigr)$\;
    $s_i^{(t)} \leftarrow s^{*} + 1$\;
  }
  \BlankLine
  $\kappa_i \leftarrow -\,\hat{\mathbf{e}}_i \cdot \mathbf{p}_i$\;
}
\end{algorithm}

% "i sees the goal": the goal is within range r of i, with no wall and no part of
% the payload on the line between them.
% All distances and directions use the shortest path across the periodic box.

% =====================================================================
% SIMPLIFICATIONS: details of the implementation (src/simulation.py)
% that were left out of the pseudocode for clarity.
% =====================================================================
%
% Curvity
%  - Code: curvity is not kappa = -(e . p). It is a piecewise-linear map of
%    e . p with three anchor values (MIN_CURVITY = -1.0 at e.p = +1,
%    MID_CURVITY = -0.25 at e.p = 0, MAX_CURVITY = +0.5 at e.p = -1).
%    With these values kappa turns positive only for headings roughly more
%    than 109 degrees from p.
%  - Curvity is recomputed every step from the current heading and the
%    stored polarity, so it also changes as the heading rotates.
%
% Update schedule
%  - Code: scores and polarities are refreshed only every
%    score_and_polarity_update_interval steps and held fixed in between.
%    The pseudocode refreshes them every step.
%  - The synchronous update is done by copying the scores once
%    (old_scores) and reading only the copy; the pseudocode expresses this
%    with the time indices (t-1) and (t).
%
% Polarity nudge (omitted entirely)
%  - Every polarity_nudge_interval steps the heading is rotated by
%    polarity_nudge_strength radians toward p (toward -p if e.p < 0).
%    This is separate from the curvity mechanism and is not shown.
%
% Isolated robots and score values
%  - Code: an isolated robot gets the integer sentinel 9999, not infinity.
%    Scores are integers (hop counts).
%  - Zero-length directions (goal exactly at the robot, or the mean
%    displacement to the best neighbours is exactly zero) return the zero
%    vector instead of a unit vector; unit(0) is left undefined above.
%
% Visibility
%  - Goal visibility: the goal must be within range r (periodic distance)
%    AND the straight line to the goal must be free of walls AND free of
%    the payload disc. The payload blocks sight of the goal, but it does
%    NOT block sight of other robots: neighbour visibility (V_i) is
%    blocked by walls only.
%  - The line-of-sight test to the goal does not wrap across the periodic
%    boundary, while the distance test does. The pseudocode treats both as
%    periodic.
%  - Neighbour wall tests use the periodic image of the neighbour.
%  - Walls are line segments or circular arcs (a curvature parameter per
%    wall); the pseudocode just says "wall".
%
% Polarity direction
%  - Polarity points toward the mean of the relative displacements to the
%    best-scoring neighbours (all weighted equally), not toward a weighted
%    or per-neighbour-normalised average. An alternative that averages unit
%    vectors exists in the code (compute_polarity_toward_minscore_ang) but
%    is not used.
%  - CLAUDE.md and docs/algorithm_notes.md describe a DIRECTEDNESS blend
%    between a Vicsek-style weighted alignment term and the gradient term.
%    The function transcribed here only contains the gradient term; the
%    blend is not shown.
%
% Efficiency machinery (no effect on the result)
%  - Neighbour search uses a cell list with cell size equal to the view
%    range and scans the 3x3 block of cells around the robot.
%  - Walls are looked up through a precomputed spatial grid.
%  - Neighbour search is two-pass (first the minimum score, then the
%    average over the minimum-score neighbours) to avoid temporary arrays.
%  - The goal test first does a cheap bounding-box rejection.
%  - The per-robot loop runs in parallel (numba prange).
%
% Left out as outside the communication layer
%  - Forces (payload, wall), the force-aligning heading update with
%    rotational noise, self-propulsion, position update and the periodic
%    wrap. Particle-particle repulsion is switched off in the code.
%  - Goal-reached detection and the end-of-run logic in runner.py.

\end{document}
